%-------------------------
% Resume in Latex
% Author : Jake Gutierrez
% Based off of: https://github.com/sb2nov/resume
% License : MIT
%------------------------

\documentclass[letterpaper,11pt]{article}

\usepackage{latexsym}
\usepackage[empty]{fullpage}
\usepackage{titlesec}
\usepackage{marvosym}
\usepackage[usenames,dvipsnames]{color}
\usepackage{verbatim}
\usepackage{enumitem}
\usepackage[hidelinks]{hyperref}
\usepackage{fancyhdr}
\usepackage[english]{babel}
\usepackage{tabularx}
\input{glyphtounicode}


%----------FONT OPTIONS----------
% sans-serif
% \usepackage[sfdefault]{FiraSans}
% \usepackage[sfdefault]{roboto}
% \usepackage[sfdefault]{noto-sans}
% \usepackage[default]{sourcesanspro}

% serif
% \usepackage{CormorantGaramond}
% \usepackage{charter}


\pagestyle{fancy}
\fancyhf{} % clear all header and footer fields
\fancyfoot{}
\renewcommand{\headrulewidth}{0pt}
\renewcommand{\footrulewidth}{0pt}

% Adjust margins
\addtolength{\oddsidemargin}{-0.5in}
\addtolength{\evensidemargin}{-0.5in}
\addtolength{\textwidth}{1in}
\addtolength{\topmargin}{-.5in}
\addtolength{\textheight}{1.0in}

\urlstyle{same}

\raggedbottom
\raggedright
\setlength{\tabcolsep}{0in}

% Sections formatting
\titleformat{\section}{
  \vspace{-4pt}\scshape\raggedright\large
}{}{0em}{}[\color{black}\titlerule \vspace{-5pt}]

% Ensure that generate pdf is machine readable/ATS parsable
\pdfgentounicode=1

%-------------------------
% Custom commands
\newcommand{\resumeItem}[1]{
  \item\small{
    {#1 \vspace{-2pt}}
  }
}

\newcommand{\resumeSubheading}[4]{
  \vspace{-2pt}\item
    \begin{tabular*}{0.97\textwidth}[t]{l@{\extracolsep{\fill}}r}
      \textbf{#1} & #2 \\
      \textit{\small#3} & \textit{\small #4} \\
    \end{tabular*}\vspace{-7pt}
}

\newcommand{\resumeSubSubheading}[2]{
    \item
    \begin{tabular*}{0.97\textwidth}{l@{\extracolsep{\fill}}r}
      \textit{\small#1} & \textit{\small #2} \\
    \end{tabular*}\vspace{-7pt}
}

\newcommand{\resumeProjectHeading}[2]{
    \item
    \begin{tabular*}{0.97\textwidth}{l@{\extracolsep{\fill}}r}
      \small#1 & #2 \\
    \end{tabular*}\vspace{-7pt}
}

\newcommand{\resumeSubItem}[1]{\resumeItem{#1}\vspace{-4pt}}

\renewcommand\labelitemii{$\vcenter{\hbox{\tiny$\bullet$}}$}

\newcommand{\resumeSubHeadingListStart}{\begin{itemize}[leftmargin=0.15in, label={}]}
\newcommand{\resumeSubHeadingListEnd}{\end{itemize}}
\newcommand{\resumeItemListStart}{\begin{itemize}}
\newcommand{\resumeItemListEnd}{\end{itemize}\vspace{-5pt}}

%-------------------------------------------
%%%%%%  RESUME STARTS HERE  %%%%%%%%%%%%%%%%%%%%%%%%%%%%


\begin{document}

%----------HEADING----------
\begin{center}
    \textbf{\Huge \scshape George Babik} \\ \vspace{1pt}
    \small 514-466-4895 $|$ \href{mailto:gbabik@student.ubc.ca}{\underline{gbabik@student.ubc.ca}} $|$ 
    \href{https://linkedin.com/in/gbabik}{\underline{linkedin.com/in/gbabik}} $|$
    \href{https://github.com/georgebabik}{\underline{github.com/georgebabik}}
\end{center}

%-----------EDUCATION-----------
\section{Education}
  \resumeSubHeadingListStart
    \resumeSubheading
      {University of British Columbia}{Vancouver, BC}
      {Bachelor of Applied Science in Electrical Engineering \textbf{(4\textsuperscript{th} Year)}}{Sep. 2022 -- May 2027}
      \resumeItemListStart
        \resumeItem{\textbf{Courses:} Digital Systems Design (CPEN 311), Computing Systems (CPEN 211), ECE Design Studio (ELEC 291)}
        \resumeItem{\textbf{Co-op:\hspace{0.3cm}} Available for 12 to 16 months beginning May 2026}
      \resumeItemListEnd
  \resumeSubHeadingListEnd

%-----------TECHNICAL SKILLS-----------
\section{Technical Skills}
\begin{itemize}[leftmargin=0.15in, label={}]
    \small{\item{
    \begin{tabbing}
        \textbf{Software Tools:} \hspace{0.2cm}\= MATLAB, ModelSim, Quartus Prime 18.1, LTspice, Git, Fusion 360, AutoCAD, LaTeX, \\
        \> Microsoft Office Suite, Xilinx Vivado, Jupyter Notebook \\[3pt]
        
        \textbf{Programming:} \> Verilog/SystemVerilog, ARM/Assembly, C, Python, Tcl \\[3pt]
        
        \textbf{Lab Equipment:} \> DE1-SoC (FPGA Board), Multimeter, Oscilloscope, Waveform Generator, Breadboard \\[3pt]

        \textbf{Languages:} \> English, Metropolitan French, Levantine Arabic
    \end{tabbing}
    }}
\end{itemize}

%-----------WORK EXPERIENCE-----------
\section{Work Experience}
  \resumeSubHeadingListStart

    \resumeSubheading
      {Wired IP ASIC Design Engineer Co-op}{January 2026 -- Present}
      {AMD Inc. (Advanced Micro Devices)}{Ottawa, ON}
      \resumeItemListStart
        \resumeItem{Supporting wired networking ASIC IP development through FPGA-based emulation, combining RTL refinement with software-driven validation workflows}
        \resumeItem{Developing software tools to generate test vectors, exercise FPGA interfaces, and capture results over high-performance PCIe links for near–real-time validation}
        \resumeItem{Collaborating with design and verification teams on simulation debugging, design reviews, and documentation for ASIC, FPGA, and Adaptive SoC platforms}
      \resumeItemListEnd

    \resumeSubheading
      {Digital Design Engineer Co-op}{September 2025 -- December 2025}
      {Blue Sky Spectroscopy}{Lethbridge, AB}
      \resumeItemListStart
        \resumeItem{Designed and automated a multichannel PID temperature control system using a Teensy 4.1 with ADC, DAC, and Lakeshore sensors for cryogenic experiments}
        \resumeItem{Implemented and evaluated PID autotuning algorithms to achieve millikelvin-level temperature stability across multiple channels}
        \resumeItem{Built and extended a Python (PySide6) GUI for real-time visualization, CSV logging, and calibration analysis of temperature and power data}
      \resumeItemListEnd

  \resumeSubHeadingListEnd

%-----------PROJECTS-----------
\section{Projects}
  \resumeSubHeadingListStart
    \resumeProjectHeading
    {{ARC4 Decryption and Brute-force Cracking Circuit}}{November 2024 -- December 2024}
    \resumeItemListStart
      \resumeItem{Developed an ARC4 symmetric stream cipher decryption circuit using a DE1-SoC FPGA board in SystemVerilog, effectively managing multiple on-chip memories and variable latencies via ready-enable interfaces}
      \resumeItem{Implemented a brute-force cracking extension capable of cycling through the entire ARC4 key space to perform successful message decryption}
      \resumeItem{Optimized brute-force performance with two parallel decryption units, accelerating the cracking process by 40\%}
    \resumeItemListEnd

    \resumeProjectHeading
    {Embedded Neural Network Accelerator}{September 2024 -- October 2024}
    \resumeItemListStart
      \resumeItem{Designed a Multi-Layer Perceptron (MLP) hardware accelerator on a DE1-SoC FPGA board in SystemVerilog, integrating an embedded Nios II CPU and off-chip SDRAM using Quartus Prime's Platform Designer}
      \resumeItem{Implemented RTL modules performing efficient matrix-vector multiplication using Q16.16 fixed-point arithmetic for accurate inference on the MNIST handwritten digit dataset}
      \resumeItem{Performed extensive debugging through JTAG interfacing and signal-level simulations, optimizing timing constraints and ensuring reliable system operation}
    \resumeItemListEnd

    \resumeProjectHeading
    {Real-time Digital Temperature Sensor and Cooling Fan}{May 2024 -- June 2024}
    \resumeItemListStart
      \resumeItem{Engineered an Arduino Uno-based temperature sensing system that activated LEDs and controlled motor speed at different temperature thresholds (below 10°C, 35-45°C, and above 45°C)}
      \resumeItem{Displayed real-time readings on an 8×5 dot matrix LCD with 1-second updates, enabling faster response to environmental changes and improving system adaptability by 20\%}
    \resumeItemListEnd
  \resumeSubHeadingListEnd
\end{document}
